\section{总结与延伸}

Switch Transformer 的核心不是“用一个 router 魔法般获得万亿模型”，而是把条件参数扩展写成一套端到端系统：router 决定路径，capacity 把动态路径约束成静态批次，load balance 提高硬件利用率，selective precision 与初始化保护训练，expert dropout 处理下游过拟合，expert parallelism 承担参数分布，实验则必须把 step、time、FLOPs 与总成本分开。

\begin{enumerate}
  \item \textbf{稀疏扩展的是可选参数池}：总参数可随 expert 数增长，每 token active experts 仍由 top-$k$ 控制；
  \item \textbf{Top-1 是 Pareto 选择}：它减少 expert FFN、通信与容量需求，但是否胜出取决于硬件和目标质量；
  \item \textbf{Router 是数值控制路径}：softmax 与 argmax 对舍入敏感，局部 float32 可以用极小成本避免路径翻转；
  \item \textbf{Load balance 是算法与系统的接口}：辅助 loss 把可学习概率变成近似规则的 expert batch；
  \item \textbf{Capacity factor 是显式成本旋钮}：更高容量减少 dropping，却扩大 buffer、通信和潜在计算；
  \item \textbf{负结果同样重要}：把 overflow token 送往次优 experts 并不必然优于跳过 expert 子层；
  \item \textbf{Scaling 必须按多个横轴验证}：step、wall-clock、FLOPs、参数存储与成本回答不同问题；
  \item \textbf{Parallelism 可以组合}：data、model 与 expert parallelism 分别切分样本、张量与条件参数；
  \item \textbf{下游相关性不等于机制证明}：upstream loss、总参数、active compute 与数据覆盖需要匹配控制；
  \item \textbf{蒸馏是有损部署转换}：应比较同规模 scratch student，并核算 teacher 训练成本；
  \item \textbf{Vision priority routing 走向自适应计算}：capacity 小于 1 时，模型主动决定哪些 patches 获得 expert 计算；
  \item \textbf{适用场景决定价值}：MoE 更偏向大规模训练与中高吞吐，不自动适合低 batch、边缘设备或极长上下文。
\end{enumerate}

\begin{importantbox}{课程作业：做一次 matched-resource MoE 审计}
实现一个小型 dense Transformer 与 top-1 MoE，匹配训练 token、active FLOPs 与总 wall-clock。至少扫描 expert 数、capacity factor 和 aux-loss 权重；报告 PPL、time-to-quality、overflow、expert entropy、all-to-all 时间、峰值显存与 per-task fine-tuning。再加入一个同参数量 dense baseline 和一个同 active-FLOPs dense baseline，回答收益来自容量、计算、正则还是系统实现。
\end{importantbox}

\begin{warningbox}{最终自检：六个不能混为一谈的概念}
\textbf{总参数}不等于\textbf{激活参数}；\textbf{FLOPs}不等于\textbf{延迟}；\textbf{负载均衡}不等于\textbf{专家相同}；\textbf{token dropping}不等于\textbf{删除 token}；\textbf{upstream correlation}不等于\textbf{参数因果作用}；\textbf{sparsity}不等于\textbf{任意 adaptive computation}。
\end{warningbox}

\subsection{拓展阅读}

建议按“历史机制—系统扩展—Switch 配方—多语言/视觉扩展”的顺序阅读原始材料：

{\small
\begin{itemize}[nosep,leftmargin=*]
  \item \href{https://doi.org/10.1162/neco.1991.3.1.79}{\textbf{Adaptive Mixtures of Local Experts}}：gating 与局部专家的早期形式化。
  \item \href{https://arxiv.org/abs/1701.06538}{\textbf{Sparsely-Gated MoE}}：大规模条件计算、路由与负载均衡的关键前身。
  \item \href{https://arxiv.org/abs/2001.08361}{\textbf{Scaling Laws for Neural Language Models}}：dense scaling 的经验背景。
  \item \href{https://arxiv.org/abs/2006.16668}{\textbf{GShard}}：自动 sharding 与大规模条件计算系统。
  \item \href{https://arxiv.org/abs/2101.03961}{\textbf{Switch Transformers}}：top-1 routing、selective precision、capacity、实验与蒸馏主来源。
  \item \href{https://arxiv.org/abs/1910.10683}{\textbf{T5}} 与 \href{https://arxiv.org/abs/2010.11934}{\textbf{mT5}}：dense text-to-text 与 101 语言基线。
  \item \href{https://arxiv.org/abs/2106.05974}{\textbf{V-MoE}}：视觉 sparse experts、priority routing 与容量小于 1 的自适应分配。
  \item \href{https://github.com/tensorflow/mesh}{\textbf{Mesh TensorFlow}}：原始工作所处的分布式实现与 sharding 代码背景。
\end{itemize}
}

\end{document}
